% Notation of the body of the paper. The figures in figures/*.tex and the
% definitions in the sections use these macros. The appendix keeps the
% TeX of the proof map (gen/*.tex, mapmacros.tex).

% Floats never precede their place in the text (flafter) and never cross a
% section boundary (\FloatBarrier at the end of each section file).
\RequirePackage{placeins}
\RequirePackage{flafter}

% ------------------------------------------------------------------ colour
% Syntax colouring of the first paper (lrlight.sty, preamble.tex): terms in
% cTm, types, session types and type environments in cTy, processes in cPr
% (black). Points, worlds, judgments and relations stay black. The keyword
% colour of the listings (kwcolor, mapmacros.tex) is the term colour.
% Lock-free parts sit on a neutral grey box (\lf, colour lfbg), which stays
% visible behind coloured text and in greyscale.
\colorlet{kwcolor}{cTm}
\colorlet{lfbg}{black!11}
\colorlet{factcolor}{cTy}
\DeclareRobustCommand{\lf}[1]{\lfadd{#1}}
\DeclareRobustCommand{\kw}[1]{\mathord{{\color{cTm}\mathsf{#1}}}}
\DeclareRobustCommand{\tyc}[1]{\mathord{{\color{cTy}\mathsf{#1}}}}
% Type-level operators and metavariables that lrlight.sty does not declare.
\DeclareRobustCommand{\tyotimes}{\mathbin{{\color{cTy}\otimes}}}
\DeclareRobustCommand{\tyoplus}{\mathbin{{\color{cTy}\oplus}}}
\DeclareRobustCommand{\tywith}{\mathbin{{\color{cTy}\&}}}
\DeclareRobustCommand{\tymultimap}{\mathbin{{\color{cTy}\multimap}}}
\DeclareRobustCommand{\tyBox}{\mathord{{\color{cTy}\Box}}}
\DeclareRobustCommand{\tyexists}{\mathord{{\color{cTy}\exists}}}
\DeclareRobustCommand{\tyforall}{\mathord{{\color{cTy}\forall}}}
\DeclareRobustCommand{\tymu}{\mathord{{\color{cTy}\mu}}}
\DeclareRobustCommand{\tyunit}{\mathord{{\color{cTy}\ast}}}
\DeclareMetaVar\cGamma{Ty}{\Gamma}
\DeclareMetaVar\cTheta{Ty}{\Theta}
\DeclareMetaVar\cPhi{Ty}{\Phi}
\DeclareMetaVar\styZ{Ty}{Z}
\DeclareMetaVar\cstr{Ty}{C}
% One macro per role. Type variables are \tvar (alpha, lrlight.sty), the
% main ordinal index is \ix and a smaller one \ixs, so alpha is never an
% ordinal. map2tex.py writes the appendix with the same macros.
\DeclareRobustCommand{\ix}{\beta}
\DeclareRobustCommand{\ixs}{\gamma}

% ------------------------------------------------------------------ terms
\DeclareRobustCommand{\tsend}{\kw{send}}
\DeclareRobustCommand{\trecv}{\kw{recv}}
\DeclareRobustCommand{\tclose}{\kw{close}}
\DeclareRobustCommand{\tselect}{\kw{select}}
\DeclareRobustCommand{\tbranch}{\kw{branch}}
\DeclareRobustCommand{\tlink}{\kw{link}}
\DeclareRobustCommand{\tspawn}{\kw{spawn}}
\DeclareRobustCommand{\tforce}{\kw{force}}
\DeclareRobustCommand{\tnew}{\kw{new}}
\DeclareRobustCommand{\tlet}{\kw{let}}
\DeclareRobustCommand{\tin}{\kw{in}}
\DeclareRobustCommand{\tcase}{\kw{case}}
\DeclareRobustCommand{\tinj}{\kw{inj}}
\DeclareRobustCommand{\tunit}{\mathord{{\color{cTm}\ast}}}
\DeclareRobustCommand{\tlam}[2]{{\color{cTm}\lambda} #1{\color{cTm}.}\,#2}
\DeclareRobustCommand{\tpair}[2]{\mathopen{{\color{cTm}\langle}} #1, #2\mathclose{{\color{cTm}\rangle}}}
\DeclareRobustCommand{\tbox}[2]{{\color{cTm}!}(#1 = #2)}
\DeclareRobustCommand{\tletp}[4]{\tlet\ \tpair{#1}{#2} = #3\ \tin\ #4}
\DeclareRobustCommand{\tletu}[2]{\tlet\ \tunit = #1\ \tin\ #2}
\DeclareRobustCommand{\tletv}[3]{\tlet\ #1 = #2\ \tin\ #3}
\DeclareRobustCommand{\tcasee}[5]{\tcase\ #1\,(#2.\,#3)\,(#4.\,#5)}

% ------------------------------------------------------------------ ends, points, worlds
\DeclareRobustCommand{\dual}[1]{\overline{#1}}
\DeclareRobustCommand{\pplus}[1]{#1^{+}}
\DeclareRobustCommand{\fld}[2]{#1\cdot #2}
\DeclareRobustCommand{\base}{\mathsf{base}}
\DeclareRobustCommand{\Point}{\mathsf{Point}}
\DeclareRobustCommand{\cons}{\mathsf{cons}}
\DeclareRobustCommand{\ranked}{\mathsf{ranked}}
\DeclareRobustCommand{\po}{\mathsf{po}}
\DeclareRobustCommand{\ordW}[1][W]{\prec_{#1}}
\DeclareRobustCommand{\wext}{\sqsubseteq}
\DeclareRobustCommand{\wjoin}{\sqcup}
\DeclareRobustCommand{\grow}[2]{\oplus_{#1} #2}
\DeclareRobustCommand{\tree}{\mathsf{tree}}
\DeclareRobustCommand{\ech}{\mathsf{chan}}
\DeclareRobustCommand{\fends}{\mathsf{fends}}

% ------------------------------------------------------------------ types
\DeclareRobustCommand{\Close}{\tyc{Close}}
\DeclareRobustCommand{\Skip}{\tyc{Skip}}
\DeclareRobustCommand{\unfold}{\mathsf{unfold}}
\DeclareRobustCommand{\sSend}[2]{{\color{cTy}!}(#1.\,#2)}
\DeclareRobustCommand{\sRecvT}[2]{{\color{cTy}?}(#1.\,#2)}
\DeclareRobustCommand{\arr}[2]{\tymultimap^{#1}_{#2}}
\DeclareRobustCommand{\ap}{\mathsf{ap}}
\DeclareRobustCommand{\up}[1]{{\vartriangle}#1}
\DeclareRobustCommand{\payload}[2]{#1[#2\cdot\bar\jmath]}
\DeclareRobustCommand{\unl}{\mathsf{unl}}
\DeclareRobustCommand{\acyc}{\mathsf{acyc}}
\DeclareRobustCommand{\conserv}{\mathsf{conserv}}
\DeclareRobustCommand{\rankable}{\mathsf{rankable}}
\DeclareRobustCommand{\tfree}{\mathsf{tfree}}
\DeclareRobustCommand{\pfree}{\mathsf{pfree}}
\DeclareRobustCommand{\pairs}{\mathsf{pairs}}
\DeclareRobustCommand{\Inst}{\mathsf{Inst}}
\DeclareRobustCommand{\idle}{\mathsf{idle}}
\DeclareRobustCommand{\chans}{\mathsf{chans}}
\DeclareRobustCommand{\ends}{\mathsf{ends}}
\DeclareRobustCommand{\tk}{\mathsf{tk}}
\DeclareRobustCommand{\contr}{\mathsf{contr}}
\DeclareRobustCommand{\cpar}[1]{\mathrel{\parallel_{#1}}}

% the typing judgment, its effect and the held values
\DeclareRobustCommand{\tyj}[5]{#1 ; #2 \vdash^{#3} #4 : #5}
\DeclareRobustCommand{\eff}[1]{\lf{\ !\ #1}}
\DeclareRobustCommand{\held}{\triangleleft}
\DeclareRobustCommand{\pj}[4]{#1 \vdash_{#2} #3 : #4}

% ------------------------------------------------------------------ relations
\DeclareRobustCommand{\Vrel}[4]{V^{#1}_{#2}\llbracket #3 \vdash #4 \rrbracket}
\DeclareRobustCommand{\Erel}[5]{E^{#1}_{#2}\llbracket #3 \vdash #4 \rrbracket^{#5}}
\DeclareRobustCommand{\Prel}[4]{P^{#1}_{#2}\llbracket #3 \vdash #4 \rrbracket}
\DeclareRobustCommand{\Fuel}[4]{F^{#1}_{#2}\llbracket #3 \vdash #4 \rrbracket}
\DeclareRobustCommand{\LFrel}[4]{L^{#1}_{#2}\llbracket #3 \vdash #4 \rrbracket}
\DeclareRobustCommand{\ready}{\mathsf{ready}}
\DeclareRobustCommand{\offered}{\mathsf{offered}}
\DeclareRobustCommand{\irred}{\mathsf{irred}}
\DeclareRobustCommand{\Obl}{\mathsf{Obl}}
\DeclareRobustCommand{\semt}{\vDash}
\DeclareRobustCommand{\semlf}{\Vvdash}
\DeclareRobustCommand{\terminated}{\mathsf{terminated}}
\DeclareRobustCommand{\dfree}{\mathsf{deadlock\text{-}free}}
\DeclareRobustCommand{\lfree}{\mathsf{lock\text{-}free}}

% ------------------------------------------------------------------ figure layout
% A judgment header in a rule figure: its name, a rule, the boxed form.
\DeclareRobustCommand{\judghead}[2]{%
  \par\noindent\makebox[\linewidth]{%
    {\small\textsc{#1}}\ \ {\color{black!35}\leaders\hrule height 0.6ex depth -0.55ex\hfill}\ \ %
    \fbox{\ensuremath{#2}}}\par\nopagebreak}
% Rules of one judgment form, centred, at one size.
\newenvironment{rules}{\begin{proofrules}}{\end{proofrules}}
\DeclareRobustCommand{\figrule}{{\color{black!35}\hrule}}
% A rule name as a link to its figure.
\DeclareRobustCommand{\ruleref}[1]{\ref{rule:#1}}
% Float placement. Two tall rule figures never share the top of one page
% with an unbreakable listing below them.
\renewcommand{\topfraction}{0.7}
\renewcommand{\textfraction}{0.2}
\renewcommand{\floatpagefraction}{0.7}
